\section{结论}
\label{sec:conclusion}

本文针对捕获前自由漂浮空间机械臂，提出了“漏斗预设性能—周期延迟零化”两段式预设时间跟踪方法，并给出了与之配套的扰动观测与切换时序。主要结论如下。

\begin{enumerate}
  \item \textbf{结构层面。}PPC 与 PDF 在语义上正交（前者约束路径、后者约束时间），但结构上不可直接叠加——时变非线性的误差变换会破坏单周期零化所依赖的线性周期时滞结构。据此按时间分段而非按结构叠加，并把交接点取在观测器的预设收敛时刻，使两种机制的激活时刻精确对齐。

  \item \textbf{双重预设。}在阶段 I 扰动正则性 \eqref{eq:cond_L}、精确模型、无饱和、超螺旋增益条件 \eqref{eq:cond_gain} 和精确 Gramian 条件下，结论为“预设时间 \(T_o\) 进带 + 再经预设时间 \(2h\) 归零”，总预设时间 \(T_{\mathrm{pre}}=T_o+2h\)。若这些条件不能核验，则只适用本文给出的残差有界结论。进带时刻来自性能函数 \eqref{eq:rho}~的构造，该函数与观测器的收敛节奏同构且在 \(T_o\) 处二阶光滑，切换不引入由 \(\rho^{(2)}\) 引起的力矩跳变。
  \item \textbf{分段观测。}\(T_o\) 前采用奇异增益的二阶动量观测器，使 \(\hat d(t)\to d(T_o)\) 且误差界与初值无关；\(T_o\) 后采用超螺旋滑模观测器，在 \eqref{eq:cond_sta}--\eqref{eq:cond_gain} 下使观测误差恒为零，为阶段 II 提供无扰动通道（定理~\ref{thm:obs1}、\ref{thm:sta}）。该结论不依赖任何额外的 Lyapunov 收敛证书，代价是 \(T_o\) 后需 \(d\) 的二阶导数有界。

  \item \textbf{阶段 I 包络保持。}漏斗界中 \(\vartheta_i<1\) 只要求观测残差\emph{有限}而与其数值无关；但实际进带深度的定量界依赖残差上界，且执行器饱和会破坏该保证（定理~\ref{thm:ppc}）。与纯 PD 阶段 I 的输入—状态界对照（推论~\ref{cor:pd_compare}）可见，PPC 的作用不是加速收敛（收敛时间由 PDF 决定），而是把观测峰值从破坏精度降级为消耗进带余量。

  \item \textbf{显式的结构瓶颈。}指出“预设包络在全程成立”并非无条件结论，给出可离线计算的严格充分判据 \(\max_iM_{h,i}B_o<\min_i\rho_{i\infty}\)（命题~\ref{prop:stage2}），它把人工延迟 \(h\)、Gramian 增益规模与误差带宽度显式耦合，是两段式结构中唯一真正的性能牺牲点。

  \item \textbf{数值验证。}在七自由度浮基机械臂（\(20\,\mathrm{N\cdot m}\) 限幅、\(T_o=5\,\mathrm{s}\)、\(h=1\,\mathrm{s}\)）上，按 \eqref{eq:ptdo}~在动量通道实现分段观测器：阶段 I 的残差按压缩时间从 \(0.31\,\mathrm{rad/s^2}\) 指数下降，阶段 II 把观测误差压到 \(10^{-4}\,\mathrm{N\cdot m}\) 量级；预设性能阶段的包络违反时长为 \(0.178\,\mathrm{s\cdot joint}\)（同包络下纯基准反馈 \(9.8\,\mathrm{s\cdot joint}\)、归一化比峰值 \(12.0\)），进带深度 \(2.1\%\!-\!90.7\%\)。因 PPC 不约束交接速度，PDF 阶段若不整形会先出现约 \(5\times10^{-3}\,\mathrm{rad}\) 的惯性外扩；按第~\ref{sec:sim_mono}~节的三步规则整形后，交接速度由 \(3.2\times10^{-2}\) 降至 \(2.5\times10^{-4}\,\mathrm{rad/s}\)，\(\lVert e\rVert_\infty\) 单调、各关节无超调（前半段逐点单调可证；后半段个别关节过零后回升 \(\le4.5\%\rho_\infty\)），最大越界比由 \(5.38\) 降至 \(0.907<1\)、包络违反为 0。稳态残差 \(4.6\times10^{-5}\,\mathrm{rad}\) 随 \(T_s\) 线性下降（\(1\,\mathrm{ms}\) 时 \(2.3\times10^{-5}\)、\(0.5\,\mathrm{ms}\) 时 \(1.2\times10^{-5}\)），与命题~\ref{prop:approx}~的最终有界结论一致。\(T_o\) 扫描表明：\(T_o=1\,\mathrm{s}\) 因包络收缩与整形窗口的双重约束而结构性失效（阶段 I 饱和 \(53.6\%\)）；\(T_o=1.5\,\mathrm{s}\) 可行但饱和达 \(29.6\%\)；\(T_o\ge2\,\mathrm{s}\) 后稳态残差与 \(T_o\) 无关（即离散地板），且单调规则下 M4\(\equiv\)M2，整定目标归结为进带深度。数值实验还给出两条工程约束：性能包络的初始宽度必须与观测器瞬态相容（第~\ref{sec:sim_infeasible}~节），否则力矩饱和会使障碍项失去意义；奇异增益只能正则化实现（\(s_{\min}\)），其与 \(T_s\) 的比值决定交接残差与步内放大的折中。
  \item \textbf{边界与降级。}在连续时间、精确模型、精确补偿与无饱和下得到精确归零；在离散实现与数值误差下，在周期递推模型成立且 \(\varepsilon_\Delta<1\) 时，退化为含 \(\varepsilon_\Delta^k\lVert z_0\rVert\) 瞬态项的有限时间界，并具有 \(\Gamma_h\bar\Delta_r/(1-\varepsilon_\Delta)\) 的最终上界（命题~\ref{prop:approx}）。同时明确记录三类不可行命题：无观测器时归零不成立且残余带可能大于误差带；超螺旋增益下界与采样/限幅上界之间存在定量三角约束 \eqref{eq:tc_tradeoff}；Gramian 病态会使 PDF 退化为等效 PD 并在 \(\varepsilon_\Delta\ge1\) 时失去有界性保证。
\end{enumerate}

后续工作包括：\textup{(1)} 把 \(T_o\) 与超螺旋增益的扫描扩展到参数摄动与初始误差变化的蒙特卡洛统计，并给出 \((\kappa_1,\kappa_2)\) 相对 \(L_i\)、以及 \(T_o\) 相对执行器预算的自整定规则；\textup{(2)} 将性能函数与增益调度结合，使 \(\rho_i\) 或 \(K_0\) 依赖构型指标，从而把本文方法与广义雅可比退化、基座反作用强度等指标联系起来；\textup{(3)} 研究时变基线 \(A_c(t)\) 下的 \(\Phi\)-加权零化增益，以消除切换跳变并容纳时变增益调度；\textup{(4)} 考虑测量噪声下的漏斗实现（噪声会使 \(\lvert\xi\rvert\) 在真值附近抖动，需要给出保包络的噪声容许条件）；\textup{(5)} 在气浮台或硬件在环平台上开展实验验证。
